\begin{abstract}
We construct an explicit counterexample to the Abhyankar--Sathaye conjecture:
a noncoordinate polynomial in four complex variables whose zero fibre is affine
three-space. Adjoining variables gives counterexamples in every ambient
dimension at least four.
\end{abstract}

\section{Introduction and statement}\label{sec:introduction}

The Abhyankar--Sathaye conjecture asks whether a polynomial
$f\in\CC[z_1,\ldots,z_n]$ satisfying
\[
  \CC[z_1,\ldots,z_n]/(f)\cong\CC^{[n-1]}
\]
must be a coordinate, that is, the first component of a polynomial automorphism
of $\Aff^n_{\CC}$.
Here $\CC^{[m]}$ denotes a polynomial ring in $m$ independent variables, and
the displayed isomorphism is required to fix $\CC$.
Geometrically, the question asks whether
every hypersurface embedding $\Aff^{n-1}_{\CC}\hookrightarrow\Aff^n_{\CC}$ can
be carried to a coordinate hyperplane by an ambient polynomial automorphism.

For $n=2$, the question is answered affirmatively by the
Abhyankar--Moh--Suzuki Theorem~\cite{AM,Suzuki}.
The higher-dimensional problem belongs to a family of questions about
polynomial maps with affine-space fibres. Sathaye~\cite[pp.~1--2]{Sathaye}
distinguished the coordinate, or epimorphism, question from the weaker question
whether an affine-space zero fibre forces all fibres to be affine space.
A coordinate polynomial defines a projection after an automorphism, so its
fibres are all isomorphic to $\Aff^{n-1}_{\CC}$; an isomorphism of the zero fibre
alone supplies no such global trivialization.
Kraft~\cite[\S4]{Kraft} surveys the embedding problem and its relation to
algebraic group actions. Among the early affirmative results,
Sathaye~\cite{Sathaye} proved rectifiability for planes defined by
polynomials linear in one of three variables in characteristic zero, and
Russell~\cite[Theorem~2.3]{Russell} extended this result to arbitrary fields.

In four variables, Kaliman, V\'en\'ereau and Zaidenberg~\cite[Theorem~4.2(a)]{KVZ}
and Maubach~\cite[Corollaries~4.2--4.3]{Maubach} established affirmative results
for equations $a(x)y+b(x,z,t)$ with $a\ne0$: an affine-three-space zero fibre
can be straightened by an ambient automorphism. Here $x,y,z,t$ are independent
variables, and the coefficient of $y$ depends on only one of them.
Further algebraic treatments appear in work of El~Kahoui, Essamaoui and
Ouali~\cite{EEO} and Ghosh, Gupta and Pal~\cite{GGPthreefold}.
Ghosh, Gupta and Pal also study linear hypersurfaces with multivariable
coefficients under additional hypotheses on the coefficients and their
divisors~\cite[Theorem~4.11]{GGP}.
Ghosh~\cite[Theorem~1.4]{Ghosh} proves an affirmative embedding theorem in
characteristic zero for
$a(x_m)b(x_1,\ldots,x_{m-1})y+f(z,t)+x_m$, where the displayed variables are
independent, $b\ne0$, and $a$ has at least two distinct roots over the algebraic
closure.
These results illustrate the role of restrictions on the defining equation;
an abstract polynomial presentation of the quotient does not itself provide
an automorphism of the ambient polynomial ring.

We give a counterexample in ambient dimension four, and hence in every
ambient dimension at least four. The three-variable case is not addressed.

The construction starts from the cusp parametrization $(u^3,-u^2)$ and
perturbs its two entries by independent multiples of $h$.
In $R=\CC[h,u,v,w]$, put
\begin{equation}\label{eq:xyz}
  \begin{aligned}
    x&=u^3+hv,\qquad y=-u^2+hw,\\
    s&=S(h,u,v,w)
       :=2u^3v+3u^4w+h(v^2-3u^2w^2)+h^2w^3.
  \end{aligned}
\end{equation}
Since $x^2+y^3$ vanishes modulo $h$, it is divisible by $h$ in $R$.
The displayed polynomial $s$ is its quotient, so $x^2+y^3=hs$.
Define $p$ and $F$ by
\begin{equation}
  p(x,y,s)=-2s^2x+3sy^2-3s^3y,
  \qquad F=h-p(x,y,s)-1. \label{eq:F}
\end{equation}
Thus $F$ is a polynomial in the four independent variables $h,u,v,w$.

\begin{theorem}\label{thm:main}
The polynomial $F$ defined in \eqref{eq:xyz}--\eqref{eq:F} satisfies
\[
  R/(F)\cong\CC^{[3]},
  \qquad
  \grad F(2,0,-1/2,1/2)=0.
\]
Consequently $F$ is not a coordinate of $R$.
\end{theorem}

The critical point $P=(2,0,-1/2,1/2)$ lies on the fibre $F=-1$.
The zero fibre is smooth, since it is isomorphic to affine three-space.

\begin{corollary}\label{cor:dimension}
For every integer $n\ge4$, there is a polynomial
$f\in\CC[z_1,\ldots,z_n]$ such that
$\CC[z_1,\ldots,z_n]/(f)\cong\CC^{[n-1]}$ and $f$ is not a coordinate.
\end{corollary}

